% Formatvorlage f�r studentische Arbeiten der Arbeitsgruppe Software Systems
% Engineering am Institut f�r Informatik der Universit�t Hildesheim
%   erstellt von Christopher Voges, 28.04.2016
%   https://sse.uni-hildesheim.de/studium-lehre/richtlinien-fuer-ausarbeitungen/vorlagen/

% Dokumentenkopf 
\documentclass[
    11pt, % Schriftgr��e
    DIV10,
    ngerman, % f�r Umlaute, Silbentrennung etc.
    a4paper, % Papierformat
    twoside, % einseitiges Dokument
    titlepage, % es wird eine Titelseite verwendet
    parskip=half, % Abstand zwischen Abs�tzen (halbe Zeile)
    headings=normal, % Gr��e der �berschriften verkleinern
    listof=totoc, % Verzeichnisse im Inhaltsverzeichnis auff�hren
    bibliography=totoc, % Literaturverzeichnis im Inhaltsverzeichnis auff�hren
    index=totoc, % Index im Inhaltsverzeichnis auff�hren
    captions=tableheading, % Beschriftung von Tabellen unterhalb ausgeben
    final % Status des Dokuments (final/draft)
]{scrreprt}

% Einbinden der Packages
% Anpassung an Sprache (deutsch)
\usepackage[ngerman]{babel}

% Umlaute
\usepackage[latin1]{inputenc}
\usepackage[T1]{fontenc}
\usepackage{textcomp} % Euro-Zeichen etc.

% Schrift
\usepackage{lmodern}
\usepackage{relsize}

% Einbinden von JPG-Grafiken erm�glichen
\usepackage[dvips,final]{graphicx}

% Erm�glichen mathematischer Symbole
\usepackage{amsmath,amsfonts}

% F�r die Definition der Zeilenabst�nde, Seitenr�nder etc.
\usepackage{setspace}
\usepackage{geometry}


% URL-Unterst�tzung
\usepackage{url}

% Abk�rzungsverzeichnis 
% Alles weitere hierzu in: "Inhalt\Abkuerzungen.tex".
\usepackage[intoc]{nomencl}
\let\abbrev\nomenclature
\renewcommand{\nomname}{Abk�rzungsverzeichnis}
\setlength{\nomlabelwidth}{.15\textwidth}

% Erm�glicht Zitate
\usepackage[square]{natbib}

% PDF-Optionen -----------------------------------------------------------------
\usepackage[
    bookmarks, % Es werden Bookmarks verwendet
    bookmarksopen=true, % Farbe von Bookmarks
    colorlinks=true, % Farbe von Verkn�pfungen
    linkcolor=black, % einfache interne Verkn�pfungen
    anchorcolor=black, % Ankertext
    citecolor=black, % Verweise auf Literaturverzeichniseintr�ge im Text
    filecolor=black, % Verkn�pfungen, die lokale Dateien �ffnen
    menucolor=black, % Acrobat-Men�punkte
    urlcolor=black, % Farbe der URLs
    plainpages=false, % zur korrekten Erstellung der Bookmarks
    pdfpagelabels, % zur korrekten Erstellung der Bookmarks
    hypertexnames=false, % zur korrekten Erstellung der Bookmarks
    linktocpage, % Seitenzahlen anstatt Text im Inhaltsverzeichnis verlinken
    pdfusetitle % Erm�glicht das Setzen der Meta-Daten des erzeugten PDFs
]{hyperref}
% \renewcommand{\theHsection}{\thepart.section.\thesection}
% \hypersetup{
%     %pdftitle={\titel \untertitel},
%     pdfauthor={\autor},
%     pdfcreator={\autor}
%     %pdfsubject={\titel \untertitel},
%    % pdfkeywords={\titel \untertitel}
% }

% Wird f�r Teile der Formatierung des Deckblatts und die Verwendung von
% Aufz�hlungen ben�tigt
\usepackage{listings}
\usepackage{xcolor} 


% fortlaufendes Durchnummerieren der Fu�noten
\usepackage{chngcntr}

% bei der Definition eigener Befehle ben�tigt
\usepackage{ifthen}

% sorgt daf�r, dass Leerzeichen hinter parameterlosen Makros nicht als Makroendezeichen interpretiert werden
\usepackage{xspace}

% Erlaubt das Anpassen der Kopf- und Fu�zeilen
\usepackage{fancyhdr}

% Erlaubt das Anpassen der �berschriften
\usepackage{titlesec}

% F�r das Erstellen eines Glossars
\usepackage[toc, automake, nonumberlist]{glossaries}

\usepackage[figure,table,lstlisting]{totalcount}

% Einbinden der Meta-Informationen
% Meta-Informationen 
% Falls Umlaute oder ein *�* vorkommen:
\usepackage[latin1]{inputenc}

% Hier k�nnen Sie Informationen zur Arbeit, sich selbst und Ihren Betreuern
% hinterlegen.
\newcommand{\titel}{Titel der Arbeit} % Name der Arbeit
\newcommand{\untertitel}{Untertitel} % Optional mit Untertitil
% Art der Arbeit ggf. zus�tzlich der Titel der Veranstaltung
\newcommand{\art}{Seminar-/Projekt-/Abschlussarbeit} 
\newcommand{\studiengang}{IMIT/WInf (BSc/MSc)} % Ihr Studiengang
\newcommand{\autor}{Max Muster} % Ihr Name
\newcommand{\email}{musterm@uni-hildesheim.de}% Ihre aktuelle und g�ltige E-Mail-Adresse
\newcommand{\matnr}{123456} % Ihre Matrikelnummer

% Die Angaben zu uns:
\newcommand{\institut}{Institut f�r Informatik}
\newcommand{\arbeitsgruppe}{Arbeitsgruppe Software Systems Engineering}
\newcommand{\erstgutachter}{Prof. Dr. Klaus Schmid, SSE}% Ihr(e) ErstgutachterIn
\newcommand{\zweitgutachter}{}% Ihr(e) ZweitgutachterIn
\newcommand{\universitaet}{Universit�t Hildesheim \textbullet  Universit�tsplatz 1 \textbullet  D-31134 Hildesheim}
\newcommand{\adresse}{\arbeitsgruppe  \textbullet  \institut \\ \universitaet}

\newcommand{\version}{Version 1.0}% Die Version der Arbeit

% Wird 'projektarbeit' auf 'true' gesetzt, wird keine Eigenst�ndigkeitserkl�rung
% erzeugt.
\newboolean{projektarbeit}
\setboolean{projektarbeit}{false}

% Wird 'final' auf 'true' gesetzt, werden folgende �nderungen vorgenommen:
% -Entfernen von Datum in der Kopf- und Versionsnummer in der Fu�zeile
% -Entfernen von Datum und Versionsnummer vom Deckblatt
% -Es werden Leerseiten f�r den doppelseitigen Druck eingef�gt
\newboolean{final}
\setboolean{final}{false}





% Erstellung der Verzeichnisse und Glossars aktivieren
\makeindex
\makenomenclature
\makeglossaries 
\glstoctrue 

% Kopf- und Fu�zeilen, Seitenr�nder etc. anpassen
% Zeilenabstand: einfach 
\newcommand{\zeilenabstandHauptteil}{1.0}
\newcommand{\zeilenabstandAnhang}{1.0}

% Seitenr�nder
\setlength{\topskip}{\ht\strutbox} % behebt Warnung von geometry
% Initiales Papierformat, wird nach dem Deckblatt teilwiese ge�ndert 
\geometry{paper=a4paper,left=3.5cm, right=2.5cm, top=2.5cm, headsep=1cm}

% Einstellen der Schriftgr��en der �berschriften
\titleformat{\chapter}[hang]{\LARGE\bfseries}{\thechapter\quad}{0pt}{}
\titleformat{\section}[hang]{\Large\bfseries}{\thesection\quad}{0pt}{}
\titleformat{\subsection}[hang]{\large\bfseries}{\thesubsection\quad}{0pt}{}
\titleformat{\subsubsection}[hang]{\large\mdseries}{\thesubsubsection\quad}{0pt}{}

% Einstellen der Abst�nde vor und nach den �berschriften
\titlespacing{\chapter}{0pt}{-3em}{0pt}
\titlespacing{\section}{0pt}{0pt}{0pt}
\titlespacing{\subsection}{0pt}{0pt}{0pt}
\titlespacing{\subsubsection}{0pt}{0pt}{0pt}
\titlespacing{\paragraph}{0pt}{0pt}{0pt}

% Definieren eines Stils f�r die Kopf- und Fu�zeilen
\fancypagestyle{mypagestyle}{%
  \fancyhf{}% Erstmal alles zur�cksetzen
  % Auf ungeraden Seiten: Kapitel�berschrift, oben rechts
  \fancyhead[OR]{\leftmark}
  \ifthenelse{\boolean{final}}{}{\fancyhead[OL]{Version vom \today~(Vor
  Abgabe entfernen)}}
  % Auf geraden Seiten: Titel der Arbeit, oben links
  \fancyhead[EL]{\titel}
  \ifthenelse{\boolean{final}}{}{\fancyhead[ER]{Version vom \today~(Vor
  Abgabe entfernen)}}
  % Jeweils die Seitenzahl
  \fancyfoot[OR]{\thepage}
  \fancyfoot[EL]{\thepage}
  % Jeweils die Version, wenn es nicht die finale Version ist, sonst nichts. 
  \ifthenelse{\boolean{final}}{}{\fancyfoot[OL]{\version~(Vor
  Abgabe entfernen)}}
  \ifthenelse{\boolean{final}}{}{\fancyfoot[ER]{\version~(Vor
  Abgabe entfernen)}}
}
\pagestyle{mypagestyle}
\renewcommand*{\chapterpagestyle}{mypagestyle} 

% So wird nur der Titel des Kapitels mit \leftmark ausgegeben.
\renewcommand{\chaptermark}[1]{\markboth{#1}{}}

%  Schriftform der Kopfzeile setzen
\renewcommand{\headfont}{\normalfont}

% Gr��e von Kopf- und Fu�zeile setzen
\renewcommand{\headrulewidth}{0.4pt}
\renewcommand{\footrulewidth}{0.4pt}

% erzeugt ein wenig mehr Platz hinter einem Punkt
\frenchspacing 

% Schusterjungen und Hurenkinder vermeiden
\clubpenalty = 10000
\widowpenalty = 10000 
\displaywidowpenalty = 10000

% Quellcode-Ausgabe formatieren
\lstset{numbers=left, numberstyle=\tiny, numbersep=5pt, breaklines=true}
\lstset{emph={square}, emphstyle=\color{red}, emph={[2]root,base}, emphstyle={[2]\color{blue}}}

% Fu�noten fortlaufend durchnummerieren
\counterwithout{footnote}{chapter}

% Einstellungen f�r Listings
\definecolor{hellgelb}{rgb}{1,1,0.9}
\definecolor{colKeys}{rgb}{0,0,1}
\definecolor{colIdentifier}{rgb}{0,0,0}
\definecolor{colComments}{rgb}{1,0,0}
\definecolor{colString}{rgb}{0,0.5,0}
\lstset{
    float=hbp,
    columns=flexible,
    tabsize=2,
    frame=single,
    extendedchars=true,
    showspaces=false,
    showstringspaces=false,
    numbers=left,
    numberstyle=\tiny,
    breaklines=true,
    breakautoindent=true,
    xleftmargin=0.6cm,
    xrightmargin=0.1cm,
    captionpos=b
}

% Eigene Definitionen f�r Silbentrennung laden
% Trennvorschl�ge im Text werden mit \" angegeben
% untrennbare W�rter und Ausnahmen von der normalen Trennung k�nnen in dieser
% Datei mittels \hyphenation definiert werden


% Eigene LaTeX-Befehle laden
% Eigene Befehle, z.B.:
\newcommand{\fett}[1]{\textbf{#1}}

% Hier beginnt das eigentliche Dokument
\begin{document}
\title{\titel}
\author{\autor}

% Setzt wie tief Abschnitte numeriert und ins Inhaltsverzeichnis
% aufgenommen werden sollen, hier: bis inkl. SubSubSection (z.B. 1.1.1.1).
\setcounter{secnumdepth}{3}
\setcounter{tocdepth}{3}

% Aus Deckblatt und Abstract Seitenzahl entfernen
\fancyfoot[OR]{}%
\fancyfoot[EL]{}%
\fancyfoot[OL]{}%
\fancyfoot[ER]{}%

% Deckblatt einbinden
% Erzeugt das Deckblatt
%   Bei zu langem Arbeitstitel m�ssen die vertikalen Abst�nde (\vspace)
%   angepasst werden, damit das Deckblatt weiterhin auf eine Seite passt.
\begin{titlepage}
\newgeometry{top=2cm,bottom=2cm,left=2cm,right=2cm}
\begin{figure}
    \begin{minipage}{0.2\textwidth}
        \begin{flushleft}    
            \includegraphics[scale=0.22]{Bilder/SSE_logo_2000_resized.pdf}
        \end{flushleft}
    \end{minipage}  
    \begin{minipage}{0.55\textwidth}
        \centering
        \hspace{0.25cm}
    \end{minipage}
    \begin{minipage}{0.2\textwidth}
        \begin{flushleft}   
            \includegraphics[scale=0.25]{Bilder/St_Uni-Logo-9-2003-eps-converted-to.pdf}
        \end{flushleft}
    \end{minipage} 
\end{figure}
\begin{center}
    \vspace*{0cm}
    \huge{\textbf{\art~im Studiengang \studiengang}}
\end{center}
\begin{center}
    \vspace{1cm}
    \Huge{\textbf{\titel}}
    \vspace{1cm}
\end{center}
\ifthenelse{\boolean{final}}{}{
    \begin{center}
        \\ \version ~vom \today        
        \\ (Vor Abgabe entfernen)
    \end{center}
}
\begin{center}
    \vspace{1cm}
    \huge{\textbf{\autor\\}}
    \vspace{1cm}
    \LARGE{\matnr\\}
    \vspace{0,5cm}
    \LARGE{\email}
\end{center}
\vspace{1cm}
\begin{center}
    \centering
    \LARGE{\textbf{Betreuer:}} \\
    \LARGE{\erstgutachter} \\
    \LARGE{\zweitgutachter} \\
\end{center}
\vspace{0.5cm}
\begin{flushright}
\end{flushright}
\begin{center}
    \small{\arbeitsgruppe  \textbullet  \institut \\ Universit�t Hildesheim
    \textbullet  Universit�tsplatz 1 \textbullet  D-31134 Hildesheim}
\end{center}
\end{titlepage}
\restoregeometry

% Setzen des Papierformats f�r den Rest des Dokuments
\newgeometry{left=3.5cm, right=2.5cm, top=2.9cm, bottom=2.9cm}

% Bei finaler Fassung: Leere Seite als R�ckseite des Deckblatts 
\ifthenelse{\boolean{final}}{\cleardoublepage}{} 

% Selbstst�ndigkeitserkl�rung einbinden
\ifthenelse{\boolean{projektarbeit}}{}{\begin{verbatim}
\end{verbatim}
\Large \textbf{Eigenst�ndigkeitserkl�rung} \normalsize
\begin{verbatim}
\end{verbatim}\setstretch{1.20}
% Bei abweichender Vorgabe bspw. durch Ihr Pr�fungsordnung, muss der
% nachfolgende Text durch den vorgegeben ersetzt werden!
\fett{Erkl�rung �ber das selbstst�ndige Verfassen von "\titel"}

Ich versichere hiermit, dass ich die vorstehende Arbeit selbstst�ndig verfasst
und keine anderen als die angegebenen Hilfsmittel benutzt habe. Die Stellen der obigen Arbeit, die anderen Werken dem Wortlaut oder dem Sinn nach entnommen wurden, habe ich in jedem Fall durch die Angabe der Quelle bzw. der Herkunft, auch der benutzten Sekund�rliteratur, als Entlehnung kenntlich gemacht. Dies gilt auch f�r Zeichnungen, Skizzen, bildliche Darstellungen sowie f�r Quellen aus dem Internet und anderen elektronischen Text- und Datensammlungen und dergleichen. Die eingereichte Arbeit ist nicht anderweitig als Pr�fungsleistung verwendet worden oder in deutscher oder einer anderen Sprache als Ver�ffentlichung erschienen. Mir ist bewusst, dass wahrheitswidrige Angaben als T�uschung behandelt werden.

Hildesheim, den~\today

\unitlength 5mm
\begin{picture}(5,4) \put(0,0) {\line(1,0){22}} \end{picture}\\
\autor
}

% Bei finaler Fassung: Leere Seite als R�ckseite der Erkl�rung
\ifthenelse{\boolean{projektarbeit}}{}{\ifthenelse{\boolean{final}}{\cleardoublepage}{}}

% Abstract einbinden
\section*{Kurzfassung}
\label{sec:Kurzfassung}
Eine kurze Zusammenfassung der Arbeit, die Interesse beim Leser wecken soll. 


\section*{Abstract}
\label{sec:Abstract}
Gerne zus�tzlich oder alternativ in Englisch.

% Bei finaler Fassung: Leere Seite als R�ckseite des Abstracts
\ifthenelse{\boolean{final}}{\cleardoublepage}{}

% Ab jetzt gibt es Seitenzahlen
\fancyfoot[OR]{\pagemark}%
\fancyfoot[EL]{\pagemark}%

% Vor dem Hauptteil werden die Seiten in gro�en r�mischen Ziffern nummeriert.
\pagenumbering{roman}

% Verzeichnisse drucken
\tableofcontents% Inhaltsverzeichnis

% Wenn Abbildungs- und Tabellenverzeichnis auf die selbe Seite sollen: 
\iftotalfigures\listoffigures\fi
\begingroup 
\let\clearpage\relax
\vspace{1cm} 
\iftotaltables\listoftables\fi
\endgroup
 
% Wenn Abbildungs- und Tabellenverzeichnis je auf einer eigenen Seite beginnen
% sollen: 
% \listoffigures % Abbildungsverzeichnis
% \listoftables % Tabellenverzeichnis

% Erstellen des Liste der Listings
\renewcommand{\lstlistlistingname}{Liste der Listings}
\iftotallstlistings\lstlistoflistings\fi

% Abk�rzungsverzeichnis einbinden
% 2 Beispiele f�r Abk�rzungen
\nomenclature{API}{Application Programming Interface}
\nomenclature{XML}{Extensible Markup Language}


% F�r korrekte �berschrift in der Kopfzeile
\clearpage\markboth{\nomname}{\nomname}

% Drucken des Abk�rzungsverzeichnises
\printnomenclature
\label{cha:Abkuerzungsverzeichnis}
 
% Arabische Seitenzahlen im Hauptteil 
\ifthenelse{\boolean{final}}{\cleardoublepage}{\clearpage}
\pagenumbering{arabic}

% Die Inhaltskapitel aus "Inhalt.tex" einbinden
\begin{spacing}{\zeilenabstandHauptteil}
% Hier k�nnen die einzelnen Kapitel inkludiert werden. 
% Die Dateien m�ssen auf .tex enden. Diese Endung muss
% beim Inkludieren aber weggelassen werden.
% Info: \include und \input unterscheiden sich im wesentlichen darin, dass bei
% \include immer eine neue Seite angefangen wird.

\chapter{Einleitung}
\label{cha:Einleitung} % Ein Label ist optional, erm�glicht aber die Referenzierung
Flie�textFlie�textFlie�textFlie�text
\section{Motivation}
\label{sec:Motivation} % Ein Label ist optional, erm�glicht aber die Referenzierung 
Flie�textFlie�textFlie�textFlie�text
\section{Ziel der Arbeit}
Flie�textFlie�textFlie�textFlie�text
\section{Aufbau der Arbeit}
Flie�text Flie�text Flie�text Flie�text Flie�text Flie�text Flie�text Flie�text
Flie�text Flie�text Flie�text Flie�text Flie�text Flie�text Flie�text Flie�text
Flie�text Flie�text Flie�text Flie�text Flie�text Flie�text Flie�text Flie�text
Flie�text Flie�text Flie�text Flie�text Flie�text Flie�text Flie�text Flie�text


\chapter{Ebenen}
In diesem Kapitel wird kurz aufgezeigt,
welche Gliederungsebenen es gibt. Hier befinden wir uns auf der Ebene eines
\fett{Kapitels}. %\fett ist ist selbst-erstelltes Kommando, dass den gleichen Effekt wie \textbf hat
\section{Abschnitt}
\subsection{Unterabschnitt}
\subsubsection{Unter-Unterabschnitt}
Soll noch tiefer gegliedert werden ist eine �nderung der Formatierungen
und/oder das Verwenden eines zus�tzlichen Paketes notwendig.

\include{Inhalt/scrivener/copy/01}
\include{Inhalt/scrivener/copy/02}
\include{Inhalt/scrivener/copy/03}
\include{Inhalt/scrivener/copy/04}
\include{Inhalt/scrivener/copy/05}
\include{Inhalt/scrivener/copy/06}
\include{Inhalt/scrivener/copy/07}
\include{Inhalt/scrivener/copy/08}
\include{Inhalt/scrivener/copy/09}
\include{Inhalt/scrivener/copy/10}
\include{Inhalt/scrivener/copy/11}
\include{Inhalt/scrivener/copy/12}
\include{Inhalt/scrivener/copy/13}
\include{Inhalt/scrivener/copy/14}
\include{Inhalt/scrivener/copy/15}
\include{Inhalt/scrivener/copy/16}
\include{Inhalt/scrivener/copy/17}
\include{Inhalt/scrivener/copy/18}
\include{Inhalt/scrivener/copy/19}
\chapter{Schlussbemerkungen}
\label{cha:Schlussbemerkungen} % Ein Label ist optional, erm�glicht aber die Referenzierung
Flie�textFlie�textFlie�textFlie�text
\section{Fazit}
Flie�textFlie�textFlie�textFlie�text
\section{Ausblick}
Flie�textFlie�textFlie�textFlie�text


\end{spacing}

% Die Inhalte des Anhangs werden analog zu den Kapiteln eingebunden.
% Wenn kein Anhang vorhanden ist, m�ssen die n�chten 7 Zeilen auskommentiert
% werden (bis '\end{spaching}').
\begin{spacing}{\zeilenabstandAnhang}
\begin{appendix}
    \chapter{Anhang}
    \label{sec:Anhang}
    % Hier k�nnen die einzelnen Anh�nge inkludiert werden. 
% Die Dateien m�ssen auf .tex enden. Diese Endung muss
% beim inkludieren aber weggelassen werden.
% Info: \include und \input unterscheiden sich im wesentlichen darin, dass bei
% \include immer eine neue Seite angefangen wird.
\section{Beispiele}
Zitat ohne Seitenangabe: \citep{BoeckleKnauberPohl+04}

Zitat mit Seitenangabe: \citep[S.1]{BoeckleKnauberPohl+04}

Referenz eines Glossareintrags: \gls{computer}

Eine normale Liste:
\begin{itemize}
\item Ein Punkt
\item Ein anderer Punkt
\end{itemize}

Eine nummerierte Liste:
\begin{enumerate}
\item Erstens
\item Zweitens
\item \ldots
\end{enumerate}

Eine Tabelle:
\begin{table}[h]
    \centering
    \captionabove{Eine Tabelle} 
    \label{tab:bspTabelle}
    \begin{tabular}{|c|c|c|} \hline
        Spalte 1 & Spalte 2 & Spalte 3 \\ \hline
        1 & 2 & 3 \\ \hline
    \end{tabular}
\end{table}

\textbf{Fetter Text}

\textit{Kursiver Text}

Eine Referenz: Siehe Tabelle \ref{tab:bspTabelle}: \nameref{tab:bspTabelle}
auf Seite \pageref{tab:bspTabelle}.

Ein Bild:
\begin{figure}[htb]
\centering
\includegraphics[width=0.2\textwidth]{Bilder/St_Uni-Logo-9-2003-eps-converted-to.pdf}
\caption{Das Logo der SUH}
\label{fig:suh-logo}
\end{figure}

Eine abgesetzte Formel:
\begin{equation}
\sum_{n=0}^{3}n=6
\end{equation}

Eine Formel im Flie�text: $2^2=4$ 

Ein Listing aus einer Datei:
\lstset{language=Java, basicstyle=\footnotesize, showstringspaces=false,tabsize=2}
\lstinputlisting[label=lst:HelloWorld,caption={HelloWorld}]{Listings/HelloWorld.java}

Ein 'on-the-fly' erstelltes Listing:
\begin{lstlisting}[[label=lst:log,caption={Beispiel eines Log-Eintrags}]] 
(Sun Sep 13 23:02:20 2009): ODBC Driver <system32>\wbemdr32.dll not present 
(Sun Sep 13 23:02:20 2009): Successfully verified WBEM OBDC adapter (incompatible version removed if it was detected).
(Sun Sep 13 23:02:20 2009): Wbemupgd.dll Registration completed.
\end{lstlisting}


 % Vor finaler Abgabe entfernen!

\end{appendix}
\end{spacing}

% Glossar einbinden
% Ein Beispiel f�r einen Glossareintrag
% Nur Eintr�ge, die mindestens ein Mal referenziert wurden 
% (z.B.: \gls{computer}), tauchen im Glossar am Ende der Arbeit auf.
 \newglossaryentry{computer} {
   name=Computer,
   description={ist ein elektronisches Ger�t, das Daten verarbeitet}
 }
\label{sec:Glossar}

% Mit diesem Befehl werden nach nicht referenierte Glossareintr�ge
% ausgegeben (Nur zum Testen einkommentieren!)
% \glsaddall

% Glossar drucken
\printglossaries
% \addcontentsline{toc}{chapter}{Glossar}

% \ifthenelse{\boolean{final}}{\cleardoublepage}{\clearpage}
\clearpage
% \addcontentsline{toc}{chapter}{Bibliography}


% Literaturverzeichnis
%   Quelldatei ist: "Bibliographie.bib".
\bibliography{Bibliographie} % Aufruf: bibtex

% Verschiedene Presets f�r den Stil der Zitate und des Literaturverzeichnisses
% \bibliographystyle{plaindin} % Nur Zahlen, deutsch
 \bibliographystyle{plain} % Nur Zahlen, englisch
%\bibliographystyle{alphadin} % Alphanumerische K�rzel, deutsch
% \bibliographystyle{alpha} % Alphanumerische K�rzel, englisch

\end{document}
